% introduction.tex — Paper 7: Chirality from h_3(O)
% ASSERT_CONVENTION: natural_units=dimensionless, jordan_product=(1/2)(ab+ba), clifford_convention=euclidean_positive, octonion_basis=fano_e1e2=e4, complex_structure=u_equals_e7, spin_representation=S10plus_boyle, pati_salam=SU4xSU2LxSU2R

\section{Introduction}
\label{sec:introduction}

The Standard Model gauge group $\SMgauge$ with its chiral fermion
representation---left-handed doublets and right-handed singlets under
$\SU{2}_L$---is the empirical target of any program that seeks to derive
particle physics from deeper principles. In a companion paper~\cite{Paper5},
we showed that \emph{self-modeling}---the requirement that a physical system
contains a faithful model of itself (a subsystem whose states separate
the states of the whole), updated through its boundary---forces
quantum mechanics: the state space of any self-modeling system is isomorphic
to $M_n(\mathbb{C})^{sa}$, the self-adjoint part of a matrix C*-algebra.

This result leaves open a structural question. Simple $M_n(\mathbb{C})$
algebras produce the gauge group $\UU{n}$ via the spectral triple
construction of Connes~\cite{CCM2007}---not the Standard Model. To obtain
the SM gauge group within the Connes--Chamseddine--Marcolli framework, the
finite algebra must be taken as a direct sum, $\mathbb{C} \oplus
\mathbb{H} \oplus M_3(\mathbb{C})$, an input chosen to match observation.
This is a \emph{structural obstruction}: self-modeling forces simple matrix
algebras, but simple matrix algebras cannot reproduce the SM within the
spectral triple approach. The obstruction is not a computational gap but a
no-go result~\cite{Barrett2007}: no Dirac operator on $M_n(\mathbb{C})$
yields the SM gauge group.

This motivates turning to a different algebraic structure: the exceptional
Jordan algebra $\hthree$, the algebra of $3 \times 3$ Hermitian matrices
over the octonions $\mathbb{O}$. The exceptional Jordan algebra occupies a
singular position in the Jordan--von Neumann--Wigner
classification~\cite{JvNW1934}: it is the unique \emph{simple} finite-dimensional
formally real Jordan algebra that is \emph{non-composable}---it cannot participate
in any composite with a nontrivial system~\cite{BGW2020}.  Under Level~IV
mathematical realism, if an algebra can form composites then composites of
it exist, making it a subsystem of a larger structure.  The universe, by
definition, is not a subsystem; therefore its algebra must be
non-composable.  This makes $\hthree$ the unique candidate for the
``universe algebra'' whose internal structure should encode
the particle content~\cite{Albert1934, Baez2002}.

The connection between $\hthree$ and the Standard Model has been explored
from several directions. Todorov and Drenska~\cite{TodorovDrenska2018}
showed that the $\Ffour = \Aut(\hthree)$ intersection with the stabilizer
of a rank-1 idempotent gives the SM gauge group. Furey~\cite{Furey2018}
showed that the Clifford algebra $\Cl{6}$, via the Witt decomposition,
produces automatic chirality for a single generation of SM fermions.
Boyle~\cite{Boyle2020} studied the complexification $\hthreeC$ and its
connection to $\Esix$. Krasnov~\cite{Krasnov2025} explored the role of pure
spinors and complex structures on $\mathbb{O}^2$.

The present paper contributes four new results to this program:
\begin{enumerate}
\item \textbf{Complexification traced to C*-observer nature.} We prove that the
complexification $\Spin{9} \to \Spin{10}$ is produced by the
\emph{combination} of the observer's C*-algebra nature (from
self-modeling~\cite{Paper5}) and the negative-spectrum Peirce operators
of $\hthree$: the continuous functional calculus of the observer's C*-algebra,
applied to these operators, yields complex coefficients, extending
the measurement algebra from
$M_{16}(\mathbb{R})$ to $M_{16}(\mathbb{C})$
(Theorem~\ref{thm:complexification}).  Neither input alone suffices---a
real-algebraic observer would have no complex functional calculus; an
operator with non-negative spectrum would have a real square root.  This
replaces the bare assumption of complexification in~\cite{Boyle2020} with
a conditional derivation from identified inputs.
This is Part~A of the argument (Sec.~\ref{sec:complexification}).
\item \textbf{One choice, two consequences.} The $\Ffour$ intersection route
(Todorov--Drenska) and the $\Cl{6}$/Pati-Salam route (Furey) are traced to
a \emph{single} algebraic input: the choice of a complex structure $u \in
S^6 \subset \im(\mathbb{O})$. Both routes produce the same gauge algebra
$\mathfrak{su}(3) \oplus \mathfrak{su}(2) \oplus \mathfrak{u}(1)$; the
$\Cl{6}$ route additionally provides the chiral representation.
\item \textbf{Symmetry reductions shown to be algebraically necessary.}  The two
symmetry-reducing steps ($E_{11}$ and $u$) are shown to be instances
of \emph{necessary symmetry reduction}: the observer must occupy
an idempotent and must complexify non-equivariantly, while all choices are
conjugate.  This reduces the conditioning inputs from three to one (Gap~A).
\item \textbf{Complete chain with explicit gaps.} No prior work connects the
self-modeling axioms through $\hthree$ to the chiral SM in a single chain
with honest gap identification. Table~\ref{tab:chain} presents the 9-link
chain from self-modeling to chiral fermions.
\end{enumerate}

\subsection{The derivation chain}
\label{subsec:chain}

Table~\ref{tab:chain} presents the complete logical chain from self-modeling
axioms to the chiral Standard Model representation. The chain consists of
nine links, each classified by source and status. One link is a genuine
conditioning input (Gap~A: Level~IV realism $+$ non-composability). Two
links are instances of \emph{necessary symmetry reduction} (B1 and
B2) in which the reduction is algebraically required and the direction is
physically irrelevant (all choices are conjugate). The remaining six links
are \emph{proved} within this work and its companions, including the
complexification step (Theorem~\ref{thm:complexification}).

The result is conditional on one input beyond the self-modeling axioms:
the identification of $\hthree$ as the universe algebra via
non-composability (Gap~A). Given this, the observer's existence forces
the Peirce decomposition (B1: $\Ffour\to\Spin{9}$, all idempotents
$\Ffour$-conjugate), the observer's C*-nature produces complexification
(Theorem~\ref{thm:complexification}), and the impossibility of
equivariant complexification (Schur's lemma) requires a choice of
$u \in S^6$ (B2: $\Gfour\to\SU{3}_C$, all $u$ $\Gfour$-conjugate).
Both symmetry reductions are algebraically necessary; both directions are gauge choices.

\begin{table*}[t]
\caption{The 9-link derivation chain from self-modeling to the chiral
Standard Model representation. Status: \textbf{Proved} = new result
derived in this work or companion papers; \textbf{Proved (standard)} =
follows from established representation theory once the prior link is
granted (the cited literature provides the computation);
\textsf{Gap} = a conditioning input not
derived from self-modeling; \textit{NSR} = necessary symmetry reduction
(direction is gauge).}
\label{tab:chain}
\begin{ruledtabular}
\begin{tabular}{clll}
Link & Statement & Source & Status \\
\midrule
L1 & Self-modeling forces $M_n(\mathbb{C})^{sa}$ (C*-algebra quantum mechanics)
   & Paper~5~\cite{Paper5} & \textbf{Proved} \\[4pt]
L2 & Non-composability of $\hthree$ identifies it as the unique ``universe algebra''
   & JvNW~\cite{JvNW1934} + BGW~\cite{BGW2020} + L4 realism
   & \textsf{Gap~A (argued)} \\[4pt]
L3 & Observer occupies rank-1 idempotent $E_{11}$; Peirce decomposition gives
     $\Vone(1) \oplus \Vhalf(16) \oplus \Vzero(10)$,
     $\Stab_{\Ffour}(E_{11}) = \Spin{9}$
   & Sec.~\ref{sec:peirce} & \textit{NSR} (B1: $\Ffour\to\Spin{9}$) \\[4pt]
L4 & C*-observer functional calculus yields complexification
     (Theorem~\ref{thm:complexification}):
     $\Vhalf \otimes_{\mathbb{R}} \mathbb{C} = \Stenplus$,
     $\Spin{9} \to \Spin{10}$
   & Sec.~\ref{sec:complexification} & \textbf{Proved} \\[4pt]
L5 & Complexification upgrades $\Ffour \to \Esix$,
     $\rep{27} \to \rep{1} \oplus \rep{10} \oplus \rep{16}$ under $\Spin{10}$
   & \cite{Baez2002,Yokota2009,Boyle2020}; Sec.~\ref{sec:spin-upgrade} & \textbf{Proved} (standard, given L4) \\[4pt]
L6 & Observer selects complex structure $u \in S^6 \subset \im(\mathbb{O})$;
     $\mathbb{O} = \mathbb{C} \oplus \mathbb{C}^3$
   & Sec.~\ref{sec:oct-split} & \textit{NSR} (B2: $\Gfour\to\SU{3}_C$) \\[4pt]
L7 & $u$ defines $\Cl{6} \subset \Cl{10}$; volume form $\omegasix$ selects
     LEFT embedding:
     $\rep{16} \to (\rep{4},\rep{2},\rep{1}) \oplus
     (\bar{\rep{4}},\rep{1},\rep{2})$
   & Sec.~\ref{sec:chirality} & \textbf{Proved} \\[4pt]
L8 & Same $u$ breaks $\Ffour \supset [\SU{3}_C \times \SU{3}_J]/\mathbb{Z}_3$;
     intersection with $\Spin{9}$ gives $\SMgauge$
   & Sec.~\ref{sec:synthesis}; $\mathrm{U}(1)$ from~\cite{TodorovDrenska2018} & \textbf{Proved} (synthesis) \\[4pt]
L9 & $\Cl{6}$/Pati-Salam route gives \textbf{same} SM gauge group as
     $\Ffour$ intersection, \textbf{plus} chiral representation
   & Sec.~\ref{sec:synthesis} & \textbf{Proved} \\
\end{tabular}
\end{ruledtabular}
\end{table*}

\subsection{Paper overview}
\label{subsec:overview}

The paper is organized as follows.

Section~\ref{sec:complexification} (Part~A) proves that the C*-algebra nature
of the observer produces complexification of the Peirce half-space $\Vhalf$,
upgrading $\Spin{9} \to \Spin{10}$ and $\Ffour \to \Esix$.  This covers
links L4--L5 of the chain.  The complexification is derived from the
observer's C*-algebra functional calculus (Theorem~\ref{thm:complexification}).

Section~\ref{sec:chirality} (Part~B) constructs the $\Cl{6}$ subalgebra
from the six directions $W = u^{\perp} \cap \im(\mathbb{O})$ defined by the
complex structure $u$. The volume form $\omegasix = \gamma_1 \cdots
\gamma_6$ selects the chiral (left-handed) embedding of the Standard Model
fermions via the Pati-Salam~\cite{PatiSalam1974} intermediate group. This
covers link~L7.

Section~\ref{sec:synthesis} presents the synthesis: the single choice of $u$
simultaneously determines both the gauge group (via the $\Ffour$ intersection
route of Todorov--Drenska~\cite{TodorovDrenska2018}) and the chiral
representation (via $\Cl{6}$). The three factors $\SU{3}_C$, $\SU{2}_L$, and
$\UU{1}_Y$ are explicitly matched across both routes. This covers links
L8--L9.

Section~\ref{sec:gaps} presents a detailed gap analysis. We identify one
genuine conditioning input (Gap~A), two instances of necessary symmetry
reduction (B1, B2), and discuss the open questions of generation
structure and the spectral action.

Section~\ref{sec:discussion} places the result in the context of the
self-modeling research program (Paper~5~\cite{Paper5}) and
discusses connections to the spectral triple approach~\cite{CCM2007},
Dubois-Violette's earlier work~\cite{DuboisViolette2016}, and the broader
literature on octonions and exceptional structures in
physics~\cite{Baez2002, Yokota2009}.
